%% CV mẫu tiếng Việt - [YOUR_NAME]
%% Dùng khi tin tuyển dụng viết bằng tiếng Việt (xem 10-vietnam-market.md).
%% Bố cục giống main_example.tex; chỉ khác tiêu đề mục, thứ tự họ tên và
%% định dạng ngày theo thói quen nhà tuyển dụng Việt Nam.
%% Latin Modern (font mặc định của moderncv khi chạy lualatex) có đủ dấu tiếng Việt.
%%
%% SETUP: chạy /setup để điền thông tin thật.
%% Compile with: lualatex main_example_vi.tex
%% (pdflatex often fails on modern MiKTeX with fontawesome5 font-expansion errors.)

\documentclass[11pt,a4paper,sans]{moderncv}
\moderncvstyle{banking}
\moderncvcolor{blue}

% Force the name and section headings to render in moderncv blue (color1).
% Default banking leaves them black: moderncvstylebanking.sty's \colorlet
% copies (not aliases) the pre-scheme accent colour, so the name colours are
% frozen before \moderncvcolor runs. Re-let them after. \namefont is the hook
% every name-style macro routes through, so this also works on moderncv 2.3.1
% (Debian/Ubuntu apt), which has no \firstnamestyle/\lastnamestyle at all.
\renewcommand*{\namefont}{\fontsize{34}{36}\bfseries\upshape}
\colorlet{firstnamecolor}{color1}
\colorlet{lastnamecolor}{color1}
\colorlet{namecolor}{color1}
\renewcommand*{\sectionstyle}[1]{{\sectionfont\color{color1}#1}}

\usepackage[utf8]{inputenc}
% pdflatex fallback only (the documented engine is lualatex, which skips this
% branch). Without T1 font encoding pdflatex builds accented letters with
% \accent, and the PDF text layer stores them decomposed - `e` + U+0300 rather
% than U+00E8 - so an ATS keyword match on "Genève" fails while the page looks
% right. moderncv 2.5 loads T1 itself under pdflatex; 2.3.1 (Debian/Ubuntu apt)
% does not. \ifpdftex comes from iftex, which every moderncv version loads.
\ifpdftex\usepackage[T1]{fontenc}\fi
% moderncv loads hyperref itself in an \AtEndPreamble hook, so \hypersetup
% must go in an \AtEndPreamble of our own: on moderncv < 2.4 a top-level
% \usepackage{hyperref} clashes with the class's own
% \RequirePackage[unicode]{hyperref}. From 2.4.0 the class passes its options
% through \PassOptionsToPackage instead, which is what removes that clash.
\AtEndPreamble{\hypersetup{
    colorlinks=true,
    linkcolor=blue,
    filecolor=magenta,
    urlcolor=blue,
    pdftitle={[YOUR_NAME] - CV},
    % Keep pdfpagemode=UseNone: this block runs after moderncv's own
    % \AtEndPreamble (moderncv.cls sets pdfpagemode there), so a FullScreen
    % value here would win and open every CV in fullscreen presentation mode.
    pdfpagemode=UseNone,
}}
\usepackage[scale=0.80]{geometry}
\usepackage{import}

% personal data
% Thứ tự tiếng Việt: \name{<Họ và tên đệm>}{<Tên>}, ví dụ \name{Nguyễn Văn}{An}
\name{[Họ và tên đệm]}{[Tên]}
% If you have no address to list, DELETE this whole line. \address{}{}{} fails
% with "There's no line here to end" on every moderncv version.
\address{[Quận/Phường, Thành phố]}{}{}
\phone[mobile]{[+84 XXX XXX XXX]}
\email{[your.email@example.com]}
\extrainfo{\href{[https://linkedin.com/in/your-profile]}{LinkedIn}, \href{[https://github.com/your-username]}{GitHub}}

% Ảnh chân dung (tuỳ chọn - phổ biến ở CV Việt Nam, nhưng công ty nước ngoài
% thường không yêu cầu). Bỏ dấu % và đặt file ảnh cạnh file .tex:
% \photo[64pt][0.4pt]{photo}

\begin{document}

\makecvtitle

% ============================================================
%     MỤC TIÊU NGHỀ NGHIỆP / GIỚI THIỆU
% ============================================================

\vspace{6pt}
\small{[3-5 dòng giới thiệu, viết riêng cho vị trí đang ứng tuyển: anh/chị là ai, mạnh nhất ở đâu, mang lại giá trị gì cho vị trí này. Ví dụ: ``Chuyên viên phân tích dữ liệu với X năm kinh nghiệm trong ngành bán lẻ, thành thạo SQL, Python và Power BI. Đã xây dựng hệ thống báo cáo doanh thu tự động cho chuỗi 50 cửa hàng. Mong muốn ứng dụng phân tích dữ liệu để hỗ trợ ra quyết định kinh doanh.'']}

% ============================================================
%     NĂNG LỰC CHÍNH
% ============================================================

\section{Năng lực chính}
\vspace{1pt}
\begin{itemize}

\item \textbf{[Nhóm kỹ năng 1]}: [Kỹ năng, công cụ cụ thể.]

\item \textbf{[Nhóm kỹ năng 2]}: [Kỹ năng, công cụ cụ thể.]

\item \textbf{[Nhóm kỹ năng 3]}: [Kỹ năng, công cụ cụ thể.]

\item \textbf{[Nhóm kỹ năng 4]}: [Chuyên môn ngành, phương pháp.]

\item \textbf{[Nhóm kỹ năng 5]}: [Phần mềm, nền tảng.]

\end{itemize}

% ============================================================
%     KINH NGHIỆM LÀM VIỆC
% ============================================================

\section{Kinh nghiệm làm việc}
\vspace{3pt}
\begin{itemize}

% --- Vị trí gần nhất ---
\item{\cventry{[MM/YYYY] - nay}{[Chức danh]}{[Công ty]}{[Thành phố]}{}{\vspace{1pt}
\begin{itemize}
    \item {[Thành tích hoặc trách nhiệm 1 - cụ thể, có con số nếu được]}
    \item {[Thành tích hoặc trách nhiệm 2]}
    \item {[Thành tích hoặc trách nhiệm 3]}
    \item {[Thành tích hoặc trách nhiệm 4]}
\end{itemize}}}

\vspace{3pt}

% --- Vị trí trước ---
\item{\cventry{[MM/YYYY] - [MM/YYYY]}{[Chức danh]}{[Công ty]}{[Thành phố]}{}{\vspace{1pt}
\begin{itemize}
    \item {[Thành tích hoặc trách nhiệm 1]}
    \item {[Thành tích hoặc trách nhiệm 2]}
    \item {[Thành tích hoặc trách nhiệm 3]}
\end{itemize}}}

\vspace{3pt}

% --- Vị trí sớm hơn ---
\item{\cventry{[MM/YYYY] - [MM/YYYY]}{[Chức danh]}{[Công ty]}{[Thành phố]}{}{\vspace{1pt}
\begin{itemize}
    \item {[Thành tích hoặc trách nhiệm 1]}
    \item {[Thành tích hoặc trách nhiệm 2]}
\end{itemize}}}

\end{itemize}

% ============================================================
%     HỌC VẤN
% ============================================================

\section{Học vấn}
\vspace{1pt}
\begin{itemize}

\item{\cventry{[YYYY-YYYY]}{[Bằng cấp] ngành [Chuyên ngành]}{[Trường]}{[Thành phố]}{}{\vspace{1pt}
[Xếp loại / GPA nếu tốt. Đề tài tốt nghiệp: ``[Tên đề tài].'']
}}

\vspace{3pt}

\item{\cventry{[YYYY-YYYY]}{[Bằng cấp] ngành [Chuyên ngành]}{[Trường]}{[Thành phố]}{}{\vspace{1pt}
[Mô tả ngắn hoặc môn học chính.]
}}

\end{itemize}

% ============================================================
%     CHỨNG CHỈ (tuỳ chọn)
% ============================================================

\section{Chứng chỉ}
\vspace{1pt}
\begin{itemize}
\item{\textbf{[Tên chứng chỉ]} -- [Đơn vị cấp] ([Năm]).}
\end{itemize}

% ============================================================
%     NGOẠI NGỮ
% ============================================================

\section{Ngoại ngữ}
\vspace{1pt}
\begin{itemize}
\item {Tiếng Việt (bản ngữ), Tiếng Anh ([IELTS/TOEIC điểm số hoặc mức độ]), [Ngôn ngữ khác] ([mức độ]).}
\end{itemize}

% ============================================================
%     GIẢI THƯỞNG (tuỳ chọn)
% ============================================================

\section{Giải thưởng}
\vspace{1pt}
\begin{itemize}
\item{\textbf{[Tên giải thưởng]} -- [Đơn vị/Sự kiện] ([Năm]).}
\end{itemize}

% ============================================================
%     NGƯỜI THAM CHIẾU
% ============================================================

\section{Người tham chiếu}
\vspace{1pt}
\begin{itemize}
\item{Sẵn sàng cung cấp khi được yêu cầu.}
\end{itemize}

\end{document}
